\documentclass[11pt,a4paper]{article}
\usepackage[utf8]{inputenc}
\usepackage[margin=1in]{geometry}
\usepackage{hyperref}
\usepackage{graphicx}
\usepackage{amsmath}
\usepackage{booktabs}

\title{RSCT Review: Causal Interpretation of Neural Network Computations with Contribution Decomposition}
\author{Swarm-It Research Discovery\\\small Automated RSCT-Based Analysis}
\date{March 09, 2026}

\begin{document}
\maketitle

\begin{abstract}
This document provides an automated review of the paper ``Causal Interpretation of Neural Network Computations with Contribution Decomposition'' from an RSCT (Representation-Solver Compatibility Testing) perspective. The paper achieved an RSCT relevance score of 30.1% and topic similarity of 47.9%.
\end{abstract}

\section{Paper Information}
\begin{itemize}
\item \textbf{Title:} Causal Interpretation of Neural Network Computations with Contribution Decomposition
\item \textbf{Authors:} Joshua Brendan Melander, Zaki Alaoui, Shenghua Liu, Surya Ganguli, Stephen A. Baccus
\item \textbf{Source:} \url{https://arxiv.org/abs/2603.06557v1}
\item \textbf{RSCT Relevance:} 30.1%
\item \textbf{Topic Match:} 47.9%
\item \textbf{Key Concepts:} representation, decomposition
\end{itemize}

\section{Original Abstract}
Understanding how neural networks transform inputs into outputs is crucial for interpreting and manipulating their behavior. Most existing approaches analyze internal representations by identifying hidden-layer activation patterns correlated with human-interpretable concepts. Here we take a direct approach to examine how hidden neurons act to drive network outputs. We introduce CODEC (Contribution Decomposition), a method that uses sparse autoencoders to decompose network behavior into sparse motifs of hidden-neuron contributions, revealing causal processes that cannot be determined by analyzing activations alone. Applying CODEC to benchmark image-classification networks, we find that contributions grow in sparsity and dimensionality across layers and, unexpectedly, that they progressively decorrelate positive and negative effects on network outputs. We further show that decomposing contributions into sparse modes enables greater control and interpretation of intermediate layers, supporting both causal manipulations of network output and human-interpretable visualizations of distinct image components that combine to drive that output. Finally, by analyzing state-of-the-art models of neural activity in the vertebrate retina, we demonstrate that CODEC uncovers combinatorial actions of model interneurons and identifies the sources of dynamic receptive fields. Overall, CODEC provides a rich and interpretable framework for understanding how nonlinear computations evolve across hierarchical layers, establishing contribution modes as an informative unit of analysis for mechanistic insights into artificial neural networks.

\section{RSCT Analysis}
This paper presents research with 30% relevance to RSCT theory.

\textbf{Abstract Summary:} Understanding how neural networks transform inputs into outputs is crucial for interpreting and manipulating their behavior. Most existing approaches analyze internal representations by identifying hidden-layer activation patterns correlated with human-interpretable concepts. Here we take a direct approach to examine how hidden neurons act to drive network outputs. We introduce CODEC (Contribution Decomposition), a method that uses sparse autoencoders to decompose network behavior into sparse mo...

\textbf{RSCT Relevance:} The work touches on concepts related to representation, decomposition, which have connections to the Representation-Solver Compatibility Testing framework.

\textbf{Further Analysis:} A detailed manual review is recommended to fully assess the implications for RSCT-based certification approaches.


\subsection*{RSCT Certification Metrics}
\begin{tabular}{ll}
\textbf{Relevance (R):} & 0.374 \\
\textbf{Spurious (S):} & 0.319 \\
\textbf{Noise (N):} & 0.307 \\
\textbf{Kappa ($\kappa$):} & 0.549 \\
\end{tabular}

The kappa score of 0.549 indicates moderate representation-solver compatibility.


\section{Relevance to Swarm-It}
This paper was identified by the Swarm-It research discovery pipeline as potentially relevant to RSCT-based AI certification. The combined score of 37.2% places it in the moderate relevance category for further investigation.

\vspace{1em}
\hrule
\vspace{0.5em}
\small{Generated by Swarm-It Research Discovery | \url{https://swarms.network} | RSCT Certified}

\end{document}
